% TeX Live 2025 XeLaTeX
\documentclass[UTF8, 12pt, a4paper, oneside, twocolumn, fontset=none]{ctexart}
\usepackage{anyfontsize}
\usepackage{geometry}
\geometry{left=1.5cm, right=1.5cm, top=1.8cm, bottom=0.9cm}

\usepackage{unicode-math, mathrsfs, xcolor, inputenc}
\definecolor{pagecolor}{HTML}{C7EDCC}
\pagecolor{pagecolor}
\setCJKmainfont[AutoFakeBold, AutoFakeSlant]{LXGW WenKai}
\setCJKmonofont[AutoFakeBold, AutoFakeSlant]{MapleMono-CN-Regular}
\setCJKsansfont[AutoFakeBold, AutoFakeSlant]{LXGW WenKai}
\setmainfont[AutoFakeBold, AutoFakeSlant]{LXGW WenKai}
\setmonofont[AutoFakeBold, AutoFakeSlant]{MapleMono-CN-Regular}
\setsansfont[AutoFakeBold, AutoFakeSlant]{LXGW WenKai}
\setmathfont{Latin Modern Math}
\setmathrm{Latin Modern Math}

\usepackage{amsmath, amsthm, cases, graphicx, subfigure, tikz, caption, subcaption, siunitx, keyval, physics2, tensor, fixdif}
\DeclareSIUnit\angstrom{Å}
\usephysicsmodule{ab, ab.braket}
\letdif*{\dr}{delta}
\letdif*{\laplacian}{Delta}
\letdif*{\nabla}{nabla}
\letdif*{\mdlgwhtsquare}{mdlgwhtsquare}
\usepackage[version=4]{mhchem}
\usepackage[flushleft]{threeparttable}
\usepackage{hyperref, bookmark}
\hypersetup{colorlinks=true, urlcolor=blue, filecolor=blue, linkcolor=blue}
\usepackage[backend=biber, sorting=none, style=authoryear, natbib=true]{biblatex}
\addbibresource{/Users/Linyan/Database/Zotero/AllBibTeX.bib}
%\addbibresource{/home/linyan/Database/Zotero/AllBibTeX.bib}

\ctexset{section={format=\bfseries\zihao{4}\flushleft}}
\renewcommand\thesubsection{\thesection.\arabic{subsection}}
\numberwithin{equation}{subsection}
\renewcommand\theequation{\thesubsection.\arabic{equation}}
\makeatletter
\renewcommand{\maketag@@@}[1]{\hbox{\m@th\normalsize\normalfont#1}}%
\makeatother
\usepackage{indentfirst}

\usepackage{listings}
\lstset{
    basicstyle      =    \zihao{-5}\ttfamily,
    numberstyle     =    \zihao{-5}\ttfamily,
    keywordstyle    =    \color{blue},
    keywordstyle    =    [2]\color{teal},
    stringstyle     =    \color{magenta},
    commentstyle    =    \color{red}\ttfamily,
    breaklines=true,     % 自动换行
    columns=fixed,       % 固定字间距
    flexiblecolumns,     % 别问为什么，加上这个
    numbers         =    left,% 行号的位置在左边
    showspaces      =    false,% 是否显示空格
    numberstyle     =    \zihao{-5}\ttfamily,% 行号的样式，小五号，tt 等宽字体
    showstringspaces=    false,
    captionpos      =    t,% 这段代码的名字所呈现的位置，t 指的是 top 上面
    frame           =    lrtb,% 显示边框
}

\lstdefinestyle{Python}{
    language        =    Python,
    basicstyle      =    \zihao{-5}\ttfamily,
    numberstyle     =    \zihao{-5}\ttfamily,
    keywordstyle    =    \color{blue},
    keywordstyle    =    [2]\color{teal},
    stringstyle     =    \color{magenta},
    commentstyle    =    \color{red}\ttfamily,
    breaklines      =    true,% 自动换行，建议不要写太长的行
    columns         =    fixed,% 如果不加这一句，字间距就不固定
    basewidth       =    0.5em,
}

\begin{document}
\section*{Homework4}

\hspace*{\fill}

\noindent\textcolor{blue}{1. Galactic Structure}

\noindent The distribution of mass in the disk of the Milky Way falls off exponentially both radially, and with perpendicular height above and below the mid-plane of the disk. The Sun is $\qty{8}{kpc}$ from the center, and $\qty{25}{pc}$ from the mid-plane. The matter density near the Sun is $\qty{4.3e-21}{kg\cdot{}m^{-3}}$.

\hspace*{\fill}

\noindent\textcolor{purple}{(a)} Using this information calculate the projected density of the Milky Way's disk as a function of radius (i.e. summing over the thickness of the disk).

\hspace*{\fill}

\noindent\textcolor{purple}{Solution}
\begin{align}
\rho\left(r, z\right)&=\rho_{0}\symup{e}^{-\frac{\left\vert{}r\right\vert}{r_{0}}}\symup{e}^{-\frac{\left\vert{}z\right\vert}{z_{0}}}.\tag{1}\\
\rho_{\odot}&=\rho_{0}\symup{e}^{-\frac{\qty{8}{kpc}}{\qty{2.5}{kpc}}}\symup{e}^{-\frac{\qty{0.025}{kpc}}{\qty{0.4}{kpc}}}=\qty{4.3e-21}{kg\cdot{}m^{-3}}.\notag\\
\rho_{0}&=\qty{1.12e-19}{kg\cdot{}m^{-3}}.\tag{2}\\
\Sigma\left(r\right)&=2\int_{0}^{\infty}\rho\left(r, z\right)\d{}z\notag\\
&=2z_{0}\rho_{0}\symup{e}^{-\frac{r}{r_{0}}}.\tag{3}
\end{align}

\hspace*{\fill}

\noindent\textcolor{purple}{(b)} If all the mass in the disk was concentrated as a thin layer of water at the mid-plane of the disk (with density $\qty{e3}{kg\cdot{}m^{-3}}$), how thick would this layer of water be?

\hspace*{\fill}

\noindent\textcolor{purple}{Solution}
\begin{align}
\Sigma\left(r\right)&=\rho_{\text{water}}h\left(r\right),\notag\\
h\left(r\right)&=2z_{0}\frac{\rho_{0}}{\rho_{\text{water}}}\symup{e}^{-\frac{r}{r_{0}}}.\tag{1}
\end{align}

\hspace*{\fill}

\noindent\textcolor{purple}{(c)} If this density was made up by stars like the Sun, what fraction of the area of the disk, as seen from a point perpendicular to the disk, would be empty?

\hspace*{\fill}

\noindent\textcolor{purple}{Solution}
\begin{align}
\Sigma_{\odot}&=\frac{M_{\odot}}{\pi{}R_{\odot}^{2}}=\qty{1.316e12}{kg\cdot{}m^{-2}}.\tag{1}\\
\eta\left(r\right)&=1-\frac{\Sigma\left(r\right)}{\Sigma_{\odot}}=1-2.1\times10^{-12}\symup{e}^{-\frac{r}{r_{0}}}.\tag{2}
\end{align}
All most everywhere is empty.

\hspace*{\fill}

\noindent\textcolor{blue}{2. Galactic Orbits}

\noindent Consider a simplified model of the Galactic Plane by assuming it to be an infinite sheet with a constant thickness and density. Show that a star displaced from the mid-plane will undergo simple harmonic motion around the mid-plane (assuming the star never leaves the region of constant density). Taking the density to be $\qty{5e6}{}$ hydrogen atoms per $\qty{}{m^{3}}$, estimate the oscillation period.

\noindent Hint: It is useful, but not necessary, to approach this using Gauss's Law of Gravity.

\hspace*{\fill}

\noindent\textcolor{purple}{Solution}

Gauss's Law of Gravity
\begin{equation}
\Phi=\oint_{S}\symbf{F}_{\symup{g}}\cdot\d{}\symbf{S}=-4\pi\symup{G}m,\tag{1}
\end{equation}
here the integral is over a closed surface and $M$ is the mass enclosed in this volume.
\begin{align}
\oint_{S}\symbf{F}_{\symup{g}1}\cdot\d{}\symbf{S}&=-2F_{\symup{g}1}S=-4\pi\symup{G}S\left(h-z\right)\rho.\tag{2}\\
\oint_{S}\symbf{F}_{\symup{g}2}\cdot\d{}\symbf{S}&=-2F_{\symup{g}2}S=-4\pi\symup{G}S\left(h+z\right)\rho.\tag{3}\\
F_{\symup{g}}&=F_{\symup{g}2}-F_{\symup{g}1}=4\pi\symup{G}\rho{}z=kz.\tag{4}
\end{align}
Note that
\begin{equation}
T=\frac{2\pi}{k},\notag
\end{equation}
then
\begin{equation}
T=\frac{2\pi}{\sqrt{4\pi\symup{G}\rho}}=\qty{2.37e15}{s}.\tag{5}
\end{equation}

%\printbibliography
\end{document}